\begin{helper}
Let \(Y_0,\ldots,Y_{r-1}\) be independent random variables taking values in a
finite set \(\Omega\), each with law \(q\).  Let \(Z\) be a real-valued function
of the product outcome \((Y_0,\ldots,Y_{r-1})\).  Suppose that changing one
coordinate while keeping all other coordinates fixed changes \(Z\) by at most
\(c\), with \(c>0\).  If \(\mu=\mathbb E Z\), then for every \(t\ge0\),
\[
\Pr(Z\ge \mu+t)\le
\exp\!\left(-\frac{2t^2}{r c^2}\right).
\]
Equivalently, when the product weight of an outcome
\(\omega=(\omega_0,\ldots,\omega_{r-1})\) is
\[
w(\omega)=\prod_{i=0}^{r-1}q(\omega_i),
\]
the total weight of all \(\omega\) with \(Z(\omega)\ge\mu+t\) is at most the
right-hand side.
\end{helper}
\begin{proof}
Order the coordinates and expose them one at a time.  For \(0\le i\le r\), let
\[
M_i=\mathbb E(Z\mid Y_0,\ldots,Y_{i-1})
\]
be the Doob martingale.  Thus \(M_0=\mu\) and \(M_r=Z\).  Fix a history of the
first \(i\) exposed coordinates and compare two possible values of \(Y_i\).
After averaging over the later independent coordinates, the two conditional
expectations differ by at most \(c\), because the two full product outcomes can
be coupled to agree in every coordinate except \(i\), and the hypothesis bounds
the resulting change in \(Z\) by \(c\).  Hence
\(\Delta_i=M_{i+1}-M_i\) has conditional mean \(0\) and, conditionally on the
past, is supported in an interval of length at most \(c\).

The needed one-step exponential estimate is Hoeffding's lemma.  In this
setting it follows directly from convexity of \(x\mapsto e^{\lambda x}\):
if a mean-zero variable \(X\) is supported in \([a,a+c]\), then the graph of
\(e^{\lambda x}\) on that interval lies below the chord joining the endpoints,
so \(\mathbb E e^{\lambda X}\) is at most the corresponding two-point endpoint
average with the same mean; maximizing that endpoint expression gives
\(\exp(\lambda^2c^2/8)\).  Applying this conditionally to each
\(\Delta_i\), for every \(\lambda>0\),
\[
\mathbb E\!\left(e^{\lambda\Delta_i}\mid Y_0,\ldots,Y_{i-1}\right)
\le \exp(\lambda^2c^2/8).
\]
Multiplying these conditional estimates along the martingale tower gives
\[
\mathbb E\exp(\lambda(Z-\mu))
=\mathbb E\exp\!\left(\lambda\sum_{i=0}^{r-1}\Delta_i\right)
\le \exp(r\lambda^2c^2/8).
\]
By Markov's inequality,
\[
\Pr(Z-\mu\ge t)
\le \exp(-\lambda t+r\lambda^2c^2/8).
\]
For \(t=0\) the desired bound is immediate.  For \(t>0\), choose
\(\lambda=4t/(r c^2)\), which minimizes the displayed quadratic exponent and
gives
\[
\Pr(Z-\mu\ge t)\le
\exp\!\left(-\frac{2t^2}{r c^2}\right).
\]
The finite-sum formulation with the weight \(w\) is exactly this product
probability written out by outcomes.
\end{proof}
